\documentclass[12pt]{article}
\usepackage[utf8]{inputenc}
\usepackage[T1]{fontenc}
\usepackage{hyperref}
\title{Tzotzil, Zinacant\'an Language Resources for K-12 Curriculum Development}
\author{Compiled by Claude Sonnet 5 (claude-sonnet-5) for Andr\'as Kornai}
\date{}
\begin{document}
\maketitle

\section{Introduction}
This report lists resources for the Zinacant\'an variety of Tzotzil (ISO 639-3 tzz; Glottocode tzot1264), spoken in Zinacant\'an, Chiapas, Mexico. The identification line in the request gives the family as ``Bookkeeping''. That looks like a data error. Tzotzil is a Mayan language, and I have not relied on the family label. The searches found mostly older scholarly reference works and archival materials. They found no complete modern K-12 curriculum, no official textbook series specific to Zinacant\'an, and no modern standalone grammar. Anything not listed below was not found.

\section{Primary Educational Resources}
\begin{itemize}
\item Tzotzil lessons and texts recorded in San Lorenzo Zinacant\'an, held in the AILLA archive (two records seen in search results; they look like language-teaching lessons, though I did not open them): \url{https://islandora-ailla.lib.utexas.edu/islandora/object/ailla%3A126555} and \url{https://islandora-ailla.lib.utexas.edu/islandora/object/ailla%3A126565}
\item A Summer Institute of Linguistics (ILV) Tzotzil primer (cartilla), San Crist\'obal de las Casas, held by El Colegio de M\'exico. It is not confirmed to be the Zinacant\'an variety: \url{https://repositorio.colmex.mx/concern/books/xg94hr58x?locale=en}
\item Sna Jtz'ibajom, a Cultural Survival project, runs community Tzotzil literacy classes, including in Zinacant\'an: \url{https://www.culturalsurvival.org/node/3833}
\end{itemize}
No official SEP Tzotzil textbooks for Zinacant\'an were found in the search results.

\section{Linguistic and Reference Materials}
\begin{itemize}
\item Laughlin, \textit{Great Tzotzil Dictionary of San Lorenzo Zinacant\'an} (Smithsonian Contributions to Anthropology 19, 1975). It is a Tzotzil--English dictionary with grammatical analysis and explanatory notes on Tzotzil grammar: \url{https://repository.si.edu/handle/10088/1360}
\item Laughlin and Haviland, \textit{Great Tzotzil Dictionary of Santo Domingo Zinacant\'an} (1988), three volumes: Tzotzil--English, English--Tzotzil and Spanish--Tzotzil. It is based on a colonial-era manuscript: \url{https://repository.si.edu/handle/10088/1367} and \url{https://repository.si.edu/handle/10088/1368}
\item World Loanword Database (WOLD) vocabulary of Zinacant\'an Tzotzil: \url{https://wold.clld.org/vocabulary/33}
\end{itemize}

\section{Religious and Community Resources}
\begin{itemize}
\item Sna Jtz'ibajom (Cultural Survival), as above: \url{https://www.culturalsurvival.org/node/3833}
\end{itemize}
Scripture translations and community organizations specific to Zinacant\'an were not found in the searches.

\section{Digital and Audio Resources}
The AILLA records listed above are the only audio or archival items found: \url{https://islandora-ailla.lib.utexas.edu/islandora/object/ailla%3A126555}. The searches did not locate any publicly archived audio collection from Laughlin's fieldwork.

\section{Curriculum Development Context}
Local schooling is contested. A news report says Tzotzil residents of Yalentay (San Joaqu\'in), in Zinacant\'an, burned free SEP textbooks: \url{https://www.eluniversal.com.mx/estados/por-segunda-ocasion-tsotziles-queman-libros-de-texto-de-la-sep-en-chiapas-buscan-ensenar-a-los-ninos-el-comunismo-dicen/}. This suggests that community involvement matters for any curriculum. The Laughlin dictionaries give a solid lexical base. Primers, readers and grammar teaching materials would still have to be developed with Zinacantec speakers, and further searching is needed before drawing conclusions about existing school materials.

\end{document}